\documentclass[11pt,letterpaper]{article}
\usepackage[utf8]{inputenc}
\usepackage[T1]{fontenc}
\usepackage[spanish,es-noshorthands,es-tabla]{babel}
\usepackage[margin=2cm,top=2.5cm,bottom=2.5cm]{geometry}
\usepackage{amsmath,amssymb}
\usepackage{enumitem}
\usepackage{fancyhdr}
\usepackage{listings}
\usepackage{xcolor}
\usepackage{booktabs}
\usepackage{array}
\raggedbottom
\setlength{\headheight}{14pt}

\pagestyle{fancy}
\fancyhf{}
\fancyhead[L]{Basic C Without Tears}
\fancyhead[C]{Guía 03 -- Tipo A}
\fancyhead[R]{Stack}
\fancyfoot[L]{Basic-C-Without-Tears}
\fancyfoot[R]{\thepage}
\renewcommand{\headrulewidth}{0.4pt}
\renewcommand{\footrulewidth}{0.4pt}

\lstdefinestyle{customc}{
  language=C,
  numbers=left,
  stepnumber=1,
  numbersep=8pt,
  numberstyle=\tiny\color{gray},
  basicstyle=\footnotesize\ttfamily,
  keywordstyle=\bfseries\color{blue!80!black},
  commentstyle=\itshape\color{green!50!black},
  stringstyle=\color{red!70!black},
  breaklines=true,
  breakatwhitespace=true,
  tabsize=4,
  showstringspaces=false,
  frame=none
}

\begin{document}
\begin{center}
    {\LARGE \textbf{Basic C Without Tears}} \\[0.25cm]
    {\Large \textbf{Guía Práctica Formativa -- Tipo A -- Guía 03}} \\[0.15cm]
    \small Fundamentos de C, traza de código, strings y matrices \\[0.1cm]
    \textbf{Autor:} Stack
\end{center}
\vspace{0.3cm}

\begin{enumerate}[label=\textbf{\arabic*.}]
    \item \textbf{Fundamentos del Lenguaje (Teoría aplicada).}
    Suponga que un \texttt{char} ocupa 1 byte y un \texttt{short} ocupa 2 bytes. Justifique las respuestas.
    \begin{enumerate}[label=(\alph*)]
        \item Para estas declaraciones, indique el valor de \texttt{sizeof(clave)}, cuántos caracteres útiles contiene la cadena y cuántos bytes suman ambos arreglos. Distinga la capacidad reservada de la longitud del texto.
        \begin{lstlisting}[style=customc,numbers=none]
char clave[7] = "LUNA";
short intentos[4];
        \end{lstlisting}
        \par\vspace{1.5cm}
        \item Indique el valor entero y el valor lógico de la expresión siguiente. Explique el efecto de \texttt{!} y de la división entera. \textquestiondown{}Se consideraría \texttt{-3} verdadero o falso si se usara solo como condición de un \texttt{if}?
        \begin{lstlisting}[style=customc,numbers=none]
!(-3) || (6 / 4 > 1)
        \end{lstlisting}
        \par\vspace{1.5cm}
        \item Indique qué imprime el fragmento. \textquestiondown{}Sería válido agregar \texttt{printf("\%d", usados);} después de cerrar el bloque? Explique el alcance de cada variable.
        \begin{lstlisting}[style=customc]
int cupos = 6;
{
    int usados = 2;
    cupos -= usados;
}
printf("%d\n", cupos);
        \end{lstlisting}
        \par\vspace{1.5cm}
    \end{enumerate}

    \newpage
    \item \textbf{Análisis y Traza de Código (Matemáticas y Control de Flujo).}
    \textquestiondown{}Qué imprime exactamente este código en pantalla? Suponga que el fragmento está dentro de \texttt{main} y que se incluyó \texttt{stdio.h}. Sin ejecutarlo ni modificarlo, respete los espacios y saltos de línea. Registre el dígito extraído, el valor restante de \texttt{n} y la suma al terminar cada vuelta, incluso cuando no se imprime nada.
    \begin{lstlisting}[style=customc]
int n = 5832, suma = 0;
while (n > 0) {
    int digito = n % 10;
    n /= 10;
    if (digito % 2 != 0) {
        continue;
    }
    suma += digito;
    printf("%d:%d\n", digito, suma);
}
printf("Resto=%d Total=%d\n", n, suma);
    \end{lstlisting}

    \newpage
    \item \textbf{Algoritmia con Strings (Pseudocódigo).}
    Un catálogo acepta una cadena solo si sus letras están en orden alfabético no decreciente: cada letra debe ser mayor o igual que la anterior. Escriba el pseudocódigo de una función que retorne \texttt{VERDADERO} o \texttt{FALSO} según esa regla, sin modificar la cadena.
    Suponga que la entrada contiene únicamente letras de \texttt{'a'} a \texttt{'z'} en ASCII. Recorra la cadena hasta \texttt{'\textbackslash 0'} y no utilice funciones de biblioteca ni ordenamiento. Las letras repetidas se permiten; la cadena vacía y una cadena de una sola letra se consideran válidas.
    \begin{center}
    \begin{tabular}{ll}
    \toprule
    \textbf{Texto} & \textbf{Resultado} \\
    \midrule
    \texttt{"abbcz"} & \texttt{VERDADERO} \\
    \texttt{"abac"} & \texttt{FALSO} \\
    \texttt{"m"} & \texttt{VERDADERO} \\
    \texttt{""} & \texttt{VERDADERO} \\
    \bottomrule
    \end{tabular}
    \end{center}

    \newpage
    \item \textbf{Algoritmia con Matrices/Arreglos (Pseudocódigo).}
    Escriba pseudocódigo que copie una matriz \texttt{origen} en otra matriz \texttt{espejo}, invirtiendo el orden de las columnas de cada fila. La primera columna debe pasar a ser la última, la segunda la penúltima, y así sucesivamente; las filas conservan su posición.
    Ambas matrices tienen \texttt{filas} por \texttt{columnas} posiciones, con dimensiones positivas y espacio ya reservado. Use bucles anidados e índices \texttt{i} y \texttt{j} desde cero, sin modificar \texttt{origen}. Muestre la matriz resultante para:
    \[
    \texttt{origen} =
    \begin{bmatrix}
    2 & 7 & 1 \\
    9 & 4 & 6
    \end{bmatrix}
    \]
    La solución debe funcionar también si hay una sola columna.
\end{enumerate}
\end{document}
